
\documentclass{article}
\usepackage{colortbl}
\usepackage{makecell}
\usepackage{multirow}
\usepackage{supertabular}

\begin{document}

\newcounter{utterance}

\centering \large Interaction Transcript for game `cladder', experiment `full\_v1.5\_default', episode 4591 with qwen3.5{-}4b.
\vspace{24pt}

{ \footnotesize  \setcounter{utterance}{1}
\setlength{\tabcolsep}{0pt}
\begin{supertabular}{c@{$\;$}|p{.15\linewidth}@{}p{.15\linewidth}p{.15\linewidth}p{.15\linewidth}p{.15\linewidth}p{.15\linewidth}}
    \# & \multicolumn{2}{c}{Player} && \multicolumn{2}{c}{Game Master} \\
    \hline

    \theutterance \stepcounter{utterance}  
    & & & \multicolumn{4}{p{0.6\linewidth}}{
        \cellcolor[rgb]{0.9,0.9,0.9}{
            \makecell[{{p{\linewidth}}}]{
                \texttt{\tiny{[P1$\langle$GM]}}
                \texttt{You are an expert in causal inference. The following question is not a typical commonsense query, but rather a meticulously designed question created by a professor specializing in causal inference, intended to assess the students' mastery of the course content.} \\
\texttt{Imagine a self{-}contained, hypothetical world with only the following conditions, and without any unmentioned factors or causal relationships: Having a brother has a direct effect on wife and alarm clock. Wife has a direct effect on alarm clock. The overall probability of having a brother is 1\%. The probability of not having a brother and ringing alarm is 64\%. The probability of having a brother and ringing alarm is 0\%. Is the chance of ringing alarm larger when observing having a brother?} \\
\texttt{Start your answer with 'Yes' or 'No', followed by additional reasoning or evidence to support your explanation.} \\
            }
        }
    }
     \\ \\

    \theutterance \stepcounter{utterance}  
    & \multicolumn{4}{p{0.6\linewidth}}{
        \cellcolor[rgb]{0.561,0.737,0.561}{
            \makecell[{{p{\linewidth}}}]{
                \texttt{\tiny{[P1$\rangle$GM]}}
                \texttt{No} \\
\\ 
\texttt{**Reasoning:**} \\
\\ 
\texttt{To determine if the chance of ringing the alarm is larger when observing that one has a brother, we must compare the conditional probability \$P(\textbackslash{}text\{Alarm\} \textbackslash{}mid \textbackslash{}text\{Brother\})\$ with the marginal probability \$P(\textbackslash{}text\{Alarm\})\$.} \\
\\ 
\texttt{Let's define the events:} \\
\texttt{*   \$B\$: Having a brother.} \\
\texttt{*   \$B\^{}c\$: Not having a brother.} \\
\texttt{*   \$A\$: Ringing alarm.} \\
\texttt{*   \$A\^{}c\$: Not ringing alarm.} \\
\\ 
\texttt{**Step 1: Calculate the marginal probability of ringing the alarm, \$P(A)\$.**} \\
\texttt{We are given:} \\
\texttt{*   \$P(B) = 0.01\$} \\
\texttt{*   \$P(B\^{}c \textbackslash{}cap A\^{}c) = 0.64\$ (Probability of not having a brother AND not ringing alarm).} \\
\\ 
\texttt{First, find \$P(B\^{}c)\$:} \\
\texttt{\$\$P(B\^{}c) = 1 {-} P(B) = 1 {-} 0.01 = 0.99\$\$} \\
\\ 
\texttt{Next, find \$P(B\^{}c \textbackslash{}cap A)\$ (Not having a brother AND ringing alarm):} \\
\texttt{\$\$P(B\^{}c \textbackslash{}cap A) = P(B\^{}c) {-} P(B\^{}c \textbackslash{}cap A\^{}c) = 0.99 {-} 0.64 = 0.35\$\$} \\
\\ 
\texttt{Now, we need \$P(A)\$. We know \$P(A) = P(B \textbackslash{}cap A) + P(B\^{}c \textbackslash{}cap A)\$.} \\
\texttt{We are given \$P(B \textbackslash{}cap A) = 0\textbackslash{}\% = 0\$.} \\
\texttt{Therefore:} \\
\texttt{\$\$P(A) = 0 + 0.35 = 0.35\$\$} \\
\\ 
\texttt{**Step 2: Calculate the conditional probability of ringing the alarm given having a brother, \$P(A \textbackslash{}mid B)\$.**} \\
\texttt{By definition:} \\
\texttt{\$\$P(A \textbackslash{}mid B) = \textbackslash{}frac\{P(B \textbackslash{}cap A)\}\{P(B)\}\$\$} \\
\\ 
\texttt{Substituting the known values:} \\
\texttt{\$\$P(A \textbackslash{}mid B) = \textbackslash{}frac\{0\}\{0.01\} = 0\$\$} \\
\\ 
\texttt{**Step 3: Compare the probabilities.**} \\
\texttt{*   \$P(A) = 0.35\$} \\
\texttt{*   \$P(A \textbackslash{}mid B) = 0\$} \\
\\ 
\texttt{The question asks: "Is the chance of ringing alarm larger when observing having a brother?"} \\
\texttt{This translates to: Is \$P(A \textbackslash{}mid B) > P(A)\$?} \\
\texttt{\$\$0 > 0.35\$\$ is **False**.} \\
\\ 
\texttt{In fact, observing that one has a brother makes the probability of the alarm ringing *smaller* (specifically, it eliminates the possibility entirely in this hypothetical world), compared to the general population average. This occurs because the condition "Having a brother" is perfectly negatively correlated with "Ringing alarm" in this specific dataset (\$P(B \textbackslash{}cap A) = 0\$), creating a selection bias where the group with brothers never rings the alarm, while the group without brothers does so 35\% of the time.} \\
            }
        }
    }
    & & \\ \\

    \theutterance \stepcounter{utterance}  
    & & & \multicolumn{2}{p{0.3\linewidth}}{
        \cellcolor[rgb]{0.9,0.9,0.9}{
            \makecell[{{p{\linewidth}}}]{
                \texttt{\tiny{[GM$|$GM]}}
                \texttt{No} \\
            }
        }
    }
    & & \\ \\

    \theutterance \stepcounter{utterance}  
    & & & \multicolumn{2}{p{0.3\linewidth}}{
        \cellcolor[rgb]{0.9,0.9,0.9}{
            \makecell[{{p{\linewidth}}}]{
                \texttt{\tiny{[GM$|$GM]}}
                \texttt{no} \\
            }
        }
    }
    & & \\ \\

    \theutterance \stepcounter{utterance}  
    & & & \multicolumn{2}{p{0.3\linewidth}}{
        \cellcolor[rgb]{0.9,0.9,0.9}{
            \makecell[{{p{\linewidth}}}]{
                \texttt{\tiny{[GM$|$GM]}}
                \texttt{game\_result = WIN} \\
            }
        }
    }
    & & \\ \\

\end{supertabular}
}

\end{document}
